%==============================================================================
%  Chapter 1 — Introduction
%
%  Aims of this chapter (from the standard "research thesis" template):
%   1. Set the empirical and policy context.
%   2. Identify the gap / puzzle.
%   3. State the research question(s) and sub-questions.
%   4. Sketch the theoretical framing.
%   5. Outline the methodological approach.
%   6. Summarise contributions.
%   7. Map the rest of the thesis.
%==============================================================================
\chapter{Introduction}
\label{ch:introduction}

% ----------------------------------------------------------------------------
\section{Background and motivation}
\label{sec:intro-background}

The European Union's economic performance is shaped by a persistent
productivity gap with the United States and China, a gap that recent
policy commentary has traced to a structural failure to scale frontier
technologies into European industrial leaders
\parencite{draghi2024future,letta2024much}\todoinline{add JRC scoreboard
or scale-up report as a third anchor}. In
the US and China, dynamic ecosystems support high firm turnover and the
emergence of new champions; in the EU, the corporate landscape remains
dominated by long-standing incumbents.

The Union has responded with a series of high-stakes initiatives that
move beyond classical market-failure correction. The European Innovation
Council (EIC), the Strategic Technologies for Europe Platform (STEP) and
the European Chips Act are not only funding programmes; they are
attempts to reform EU industrial policy by combining incremental support
for established sectors with direct backing for breakthrough
technologies. This reformist turn carries a tension that has been read
variously as a contradiction or as a deliberate re-tooling of public
capacity \parencite{mazzucato2018mission}: the Union is catalysing
radical, high-uncertainty growth on a mandate historically built around
market neutrality, competition policy and regional redistribution. The
thesis takes no position on which reading is correct, and asks instead
how that tension is handled in the text of two of the EU's most
prominent reformist instruments, Pathfinder and STEP Scale Up.

% ----------------------------------------------------------------------------
\section{Problem statement and research gap}
\label{sec:intro-problem}
Pathfinder and STEP Scale Up share an institutional home and a single
annual Work Programme, yet they appear to operate by different rules.
Pathfinder funds early-stage research consortia through non-dilutive
grants at TRL~1--4; STEP Scale Up channels equity-style catalytic
investment into companies seeking industrial expansion. Both
instruments are nonetheless presented in the Work Programme as a
coherent EIC offer, with the same language of breakthrough, autonomy
and impact running across them. The empirical puzzle is how a single
agency, working from one text, holds these two postures together: an
exploratory grant logic at one end of the pipeline and an equity-based
scale-up logic at the other, without visible reconciliation of the
assumptions each carries about applicants, evidence and success.

The puzzle invites an institutional reading because the work the Work
Programme is doing is not only allocative but constitutive. The text
defines who counts as an EIC applicant, what counts as a credible
innovation claim, and what counts as a successful outcome. Institutional
theory provides a vocabulary for that kind of textual constitution, and
the constructs the thesis mobilises are introduced in
\cref{sec:intro-theory} and developed in \cref{ch:background}. The
existing literature on EU innovation finance describes the
blended-finance turn extensively, but almost uniformly in a
policy-advocacy register: blended and mission-oriented public
investment is presented as the answer to the European scale-up gap
\parencite{mazzucato2018mission}. The reading offered here moves in a
different direction. It treats Pathfinder and STEP not as solutions to
be advocated but as institutional constructs to be interrogated, and
asks what assumptions and tensions their textual design carries. The
gap is therefore twofold. Empirically, the author is not aware of a
comparative qualitative reading of Pathfinder and STEP as co-existing
instruments inside the same Work Programme that works in this critical
register. Theoretically, institutional-logics scholarship has not yet
engaged seriously with the kind of public-investment vehicle the EIC
has become.

% ----------------------------------------------------------------------------
\section{Origin of the inquiry}
\label{sec:intro-origin}

This thesis sits within a larger observation about how European
companies are financed. The observation does not come from the
literature; it comes from the researcher's experience in venture
building, investing and consulting, and from patterns visible
from inside that work. On that prior view, the abundance of
public, non-dilutive funding in Europe shifts the stage at which
private capital enters: private rounds arrive later in a
company's life but at similar absolute sizes, leaving firms
under-capitalised at the stage at which they are then financed.
That broader view is a claim about firm-level financing
trajectories and cannot be settled with policy text alone. What
the text can show is how the public instruments at two points
relevant to that view are designed: Pathfinder at the early
stage at which the prior view holds the dynamic to begin, and
STEP at the scale-up stage to which public capital has more
recently been extended. This thesis examines those two design
moves, as an initial probe of whether the architecture exhibits
textual features consistent with the mechanism the prior view
posits. The firm-level question (whether the dynamic actually
operates as described) requires firm-level data and a different
research design, and is left for a later study.

% ----------------------------------------------------------------------------
\section{Research question and sub-questions}
\label{sec:intro-rq}

\noindent\textbf{Main research question.} \\
\textit{How are investment logics constructed, legitimised and contested
within the EIC instruments Pathfinder and STEP, and how can these dynamics
be understood through the lens of institutional theory?}

\medskip
\noindent\textbf{Sub-questions.} The sub-questions are framed to match
the documentary scope of the study: they ask what the Work Programme
text \emph{inscribes}, \emph{constructs} and \emph{prescribes}, rather
than how individual actors enact those prescriptions in practice.
\begin{enumerate}[label=\textbf{SQ\arabic*.},leftmargin=*]
    \item Which institutional logics are inscribed in the Pathfinder and
          STEP texts of the EIC Work Programme 2026, and how are they
          distributed across the instruments' architecture
          (eligibility, evaluation, governance, impact)?
    \item How does the Work Programme construct legitimacy for the
          co-presence of these logics, and through what rhetorical and
          structural devices?
    \item How does STEP reconfigure the relationship between state and
          market logic relative to Pathfinder?
    \item What does the comparison between Pathfinder and STEP reveal
          about the model of public investment the EU is
          institutionalising?
\end{enumerate}

% ----------------------------------------------------------------------------
\section{Theoretical positioning}
\label{sec:intro-theory}
The thesis reads the Work Programme through institutional theory, with
institutional logics in the Friedland--Alford and Thornton tradition
as the primary lens \parencite{friedland1991bringing,
thornton2012institutional}. Four logics are foregrounded: science,
market, state and profession. Around this primary lens, three further
constructs do supporting work. DiMaggio and Powell's isomorphism
\parencite{dimaggio1983iron} helps explain why the EIC's design echoes
external referents such as DARPA and Temasek. Suchman's legitimacy
framework names the moves by which the text justifies the co-presence
of otherwise divergent logics. Scott's three pillars
\parencite{scott2014institutions} provide a coarser background grammar
for how each logic is held in place. The institutional lens is in turn
complemented by the mission-oriented and developmental-state
literatures \parencite{mazzucato2018mission}, which supply the
comparator vocabulary against which STEP's investment mechanics will
later be read. The framework is set out in full in \cref{ch:background}.

% ----------------------------------------------------------------------------
\section{Methodological approach}
\label{sec:intro-method}
The empirical material is documentary. The thesis works from the
Pathfinder and STEP Scale Up sections of the EIC Work Programme 2026
\parencite{ec2024eicwp}, supplemented by the regulatory instruments
that establish the two programmes \parencite{eu2024step}. No interviews
were conducted, and the analytical claims are restricted to what the
text inscribes.

The analysis runs as three sequential passes over the corpus. A first,
deductive readthrough codes each substantive passage against a codebook
derived from institutional theory, producing numbered entries with
codings, quoted evidence and open flags. A second, inductive pass
resolves those flags and develops pattern-level observations across
entries. A third pass synthesises the patterns into the findings
reported in \cref{ch:results}. The full procedure, codebooks and audit
trail are documented in \cref{ch:methodology} and
\cref{app:codebook-deductive,app:codebook-inductive,app:audit-trail}.

% ----------------------------------------------------------------------------
\section{Scope and delimitations}
\label{sec:intro-scope}
Beyond the instrument scope already set out, the thesis is bounded in
two further ways.

First, non-EU public-investment vehicles appear only as comparators.
DARPA, Temasek and the Chinese Government Guidance Funds are referenced
where the Work Programme itself invokes them, or where the secondary
literature uses them to characterise mission-oriented design. They are
not treated as cases in their own right and are not subjected to the
same close reading as the EIC text.

Second, the analysis is textual and qualitative. It makes no claim
about the quantitative impact of either instrument, and it does not
examine how Programme Managers, evaluators, applicants or investee
firms enact the textual prescriptions in practice. The reading
captures what the EIC text \emph{commits} the instruments to, not what
actors then do with those commitments. Firm-level questions are left
for a later study (see \cref{sec:intro-origin}); the reflexive grounds
for the restriction are set out in \cref{sec:meth-reflexivity}.

% ----------------------------------------------------------------------------
\section{Contributions}
\label{sec:intro-contribution}
\begin{itemize}
    \item \textbf{Empirical contribution.} The thesis provides, to the
          author's knowledge, the first comparative qualitative reading
          of Pathfinder and STEP Scale Up as co-inscribed in the EIC
          Work Programme 2026 that works in a critical institutional
          register rather than in the policy-advocacy register dominant
          in the wider literature on EU blended finance. The reading is
          anchored in a numbered, traceable corpus of coded entries,
          with the audit trail reproduced in \cref{app:audit-trail}.
    \item \textbf{Practical contribution.} The reading identifies the
          specific textual mechanisms by which the Work Programme
          manages logic conflict. These mechanisms are the most
          immediate object the successor framework programme (FP10) can
          rewrite if it wants its own scale-up vehicle to handle the
          same tensions differently.
\end{itemize}

% ----------------------------------------------------------------------------
\section{Structure of the thesis}
\label{sec:intro-structure}
The remainder of the thesis is organised as follows.
Chapter~\ref{ch:background} reviews the literature on institutional theory
and on EU innovation funding, situating Pathfinder and STEP within both
streams.
Chapter~\ref{ch:methodology} describes the research design, case selection,
data sources, coding strategy and quality criteria.
Chapter~\ref{ch:results} presents the empirical findings as a logic
profile across the two instruments, four cross-cutting patterns, a
cluster of STEP-specific dynamics, the role of TRL as a
cognitive-pillar device, and the mechanisms by which sustainability
enters the Challenge chapters.
Chapter~\ref{ch:discussion} interprets the findings against the theoretical
framework and engages with rival explanations.
Chapter~\ref{ch:conclusion} summarises the contributions, acknowledges
limitations, and proposes avenues for future research.
